\section{Introduction}
\label{sec:introduction}
Begin with a concrete workload or failure mode, identify the responsible
component, and explain the observable consequence. Connect a learned
model's behavior to the software and runtime that execute it.

\paragraph{Problem and gap}
Specify what existing validation observes and misses. Preserving a reference
implementation's property through a transformation is a different question
from establishing that the reference was correct. Keep these claims separate.

\paragraph{Design insight}
Describe the interface and oracle used to compare executions. Identify
where a policy is enforced and which components it trusts. Protection
principles provide context \citep{saltzer1975}; the paper must still supply
a concrete model and supporting evidence.

The sample structure supports three kinds of contribution:
\begin{itemize}
  \item A precisely scoped problem or threat model.
  \item A system, analysis, or defense with an explicit mechanism.
  \item An evaluation of effectiveness, overhead, and failure cases.
\end{itemize}

\section{Scope and Threat Model}
\label{sec:threat}
\paragraph{Assets and boundaries}
Identify the protected property and trust boundaries. State whether the
model, compiler, dependencies, runtime, and host are trusted. Distinguish
data flow from control and privilege in diagrams.

\paragraph{Adversarial control}
List what an adversary can read, modify, or query, with the available budget
and feedback. For non-adversarial work, specify faults and workloads instead.

\paragraph{Success and exclusions}
Define success as an observable outcome. Explain exclusions such as
physical access or privileged host compromise. Exclusions narrow the
claim; they do not establish safety.
